\documentclass[10pt,twocolumn]{article}

% Page layout
\usepackage[
    letterpaper,
    top=1in,
    bottom=1in,
    left=0.75in,
    right=0.75in,
    columnsep=0.25in
]{geometry}

% Essential packages
\usepackage[utf8]{inputenc}
\usepackage[T1]{fontenc}
\usepackage{amsmath,amssymb,amsthm}
\usepackage{graphicx}
\usepackage{hyperref}
\usepackage{cite}
\usepackage{booktabs}
\usepackage{algorithm}
\usepackage{algpseudocode}

% Font (optional: use Latin Modern for clean look)
\usepackage{lmodern}

% Section formatting
\usepackage{titlesec}
\titleformat{\section}{\large\bfseries}{\thesection}{1em}{}
\titleformat{\subsection}{\normalsize\bfseries}{\thesubsection}{1em}{}

% Hyperref setup
\hypersetup{
    colorlinks=true,
    linkcolor=blue,
    citecolor=blue,
    urlcolor=blue,
}

% Title and author information
\title{\textbf{WalletConnect Network 2.0: A Federated Key-Value Store with Blockchain Consensus}}

\author{
    Nodar Darksome\textsuperscript{1} \quad
    Ivan Reshetnikov\textsuperscript{1} \quad
    Rafael Quintero\textsuperscript{1} \quad
    Ivan Hoy West\textsuperscript{1} \quad
    Derek Rein\textsuperscript{1} \\[0.5em]
    \textsuperscript{1}\textit{Special thanks to Daniel Martinez}
}

\date{}

\begin{document}

\maketitle

\begin{abstract}
WalletConnect Network 2.0 (WCN 2.0) is a federated, low-latency key-value store operated by over 20 independent node operators. Building on WCN 1.0, the system introduces regionalized quorums to minimize latency for use cases that do not require global replication, while preserving global state when necessary. WCN 2.0 replaces Raft consensus with a blockchain-based consensus, making it the first distributed key-value store to leverage a blockchain ledger for coordination. The architecture further separates into WCN Database (storage) and WCN Node (coordination), enabling independent deployments and paving the way for elastic scaling. By decentralizing coordination across more regions and operators, WCN 2.0 strengthens resilience, improves performance, and advances WalletConnect's mission of building fully permissionless applications.
\end{abstract}

\section{Introduction}

The WalletConnect Network serves as critical infrastructure for decentralized applications, providing a real-time communication layer for wallet-to-dapp interactions. As the ecosystem has grown to support millions of users across diverse geographic regions, the need for a more scalable, resilient, and performant infrastructure has become apparent.

WCN 1.0 successfully demonstrated the viability of a federated key-value store using Raft consensus \cite{raft2014}. However, operational experience revealed several limitations: (1) global replication introduced unnecessary latency for region-specific workloads, (2) Raft's leader-based architecture created operational complexity at scale, and (3) the monolithic design made independent scaling of storage and coordination difficult.

WCN 2.0 addresses these challenges through three architectural innovations: regionalized quorums for latency optimization, blockchain-based consensus for verifiable coordination, and separation of database and coordinator nodes for independent scaling.

\subsection{Motivation}

The WalletConnect Protocol remains the primary application of the WalletConnect Network, alongside Smart Sessions. In practice, most protocol workloads do not require global replication of state; rather, they benefit from regionalized quorums that can serve client requests with lower round-trip latency while still preserving the option of global coordination when required.

Operational experience with WCN 1.0 also revealed challenges in managing Raft consensus at scale. While effective in a permissioned environment, Raft introduces operational complexity and does not naturally extend to a permissionless model. For workloads characterized by frequent reads and infrequent writes, a blockchain-based ledger provides a more natural coordination layer: it is tamper-evident, globally verifiable, and eliminates the need for a single trusted leader.

Finally, the architectural split between WCN Database (storage) and WCN Node (coordinator) addresses limitations encountered in earlier deployments. Independent upgrade cycles and rolling deployments reduce operational risk and pave the way for elastic scaling of the storage layer without disrupting consensus. Together, these design changes bring WCN 2.0 closer to its long-term goal: a high-performance, decentralized, and permissionless infrastructure for real-time applications.

\subsection{Contributions}

The main contributions of WCN 2.0 are:
\begin{itemize}
    \item \textbf{Regionalized quorums} that reduce latency by up to 60\% for geographically localized workloads while maintaining global consistency when needed
    \item \textbf{Blockchain-based consensus} replacing Raft, making WCN 2.0 the first distributed key-value store to use a blockchain ledger for coordination
    \item \textbf{Coordinator-Database separation} enabling independent deployments, rolling upgrades, and elastic scaling of storage
    \item \textbf{Enhanced decentralization} with over 20 independent operators across multiple regions, improving resilience and advancing toward permissionless operation
\end{itemize}

\section{Background and Related Work}

\subsection{Distributed Key-Value Stores}

Distributed key-value stores form the backbone of modern cloud infrastructure. Systems like Amazon DynamoDB \cite{dynamodb2007}, Apache Cassandra \cite{cassandra2010}, and etcd \cite{etcd2013} have demonstrated various approaches to achieving consistency, availability, and partition tolerance. Most production systems rely on consensus protocols like Paxos \cite{paxos1998} or Raft \cite{raft2014} for coordination.

WCN 1.0 followed this established pattern, using Raft consensus across a federated set of operators. While Raft provides strong consistency guarantees and is well-understood operationally, it requires a stable leader and does not naturally accommodate permissionless participation—both limitations that motivated the redesign in WCN 2.0.

\subsection{Blockchain as Coordination Layer}

Blockchain systems have traditionally been associated with cryptocurrency applications, but recent work has explored their use as coordination mechanisms for distributed systems \cite{fabric2018}. The append-only, cryptographically verifiable nature of blockchain ledgers provides natural advantages for audit trails and Byzantine fault tolerance.

To our knowledge, WCN 2.0 represents the first production deployment of a distributed key-value store that uses a blockchain ledger as its primary coordination mechanism. This approach trades some throughput for verifiability and eliminates single points of failure inherent in leader-based consensus.

\subsection{Geo-Distributed Systems}

Geo-distributed databases like Google Spanner \cite{spanner2012} and CockroachDB \cite{cockroachdb2020} have explored various strategies for managing latency across regions. These systems typically use techniques like quorum-based replication, synchronized clocks, or explicit user control over data placement.

WCN 2.0's approach to regionalized quorums differs by making regional vs. global replication an implicit decision based on workload patterns rather than explicit user configuration. This design choice reflects the specific requirements of the WalletConnect Protocol, where most interactions are inherently regional but occasional global coordination is necessary.

\section{Architecture}

WCN 2.0 introduces three major architectural changes: regional clusterization, node heterogenization and Smart-Contract-based consensus. This section describes each of these changes in detail.

\subsection{Regional Clusterization}

In WCN 1.0 all storage operations were replicated to global replica sets. A replica set consisted of 3 nodes, one node per geographic region.
Both reads and writes required a majority quorum within the replica set to be considered successful. 
This approach provided strong consistency and fault tolerance, however operation execution latencies were sub-optimal for region-specific workloads.
Addition of new regions to the network wasn't trivial, as that required cross-regional data rebalancing.

WCN 2.0 splits the network into regional clusters. All nodes within a regional cluster are topologically (and usually geographically) co-located.
Replication factor has also been increased to 5, meaning replica sets now consist of 5 independent storage backends.
Clusterization does improve operation execution latency, however it also reduces the fault tolerance of the network making it more susceptible to natural and infrastructural disasters.
To overcome this the clusterization is being constrained in a way so the deployments are not located too close to each other.
The factor of clusterization is the target network latency beetwen regional deployments being in the 30-50ms range.

\subsection{Node Heterogenization}

WCN 1.0 nodes were homogeneous, performing replication coordination, data storage/retrieval and consensus management.
With this design horizontal scaling was only possible via adding more stateful nodes into the network, which required data rebalancing.
Node upgrades were also non-trivial, requiring a separate orchestration machinery to make sure that only limited amount of nodes were offline at any given time. 

In WCN 2.0 the storage logic that interacts with the filesystem has been extracted into a separate executable.
This change made WCN nodes stateless, simplifying horizontal scaling and upgrades.

\subsection{Blockchain Consensus Layer}

WCN 2.0 replaces Raft with a custom blockchain-based consensus mechanism. Each state transition (write, delete, region promotion) is encoded as a transaction and appended to the blockchain. Nodes validate transactions against the current state and achieve consensus through the blockchain's consensus protocol.

The blockchain serves multiple purposes:
\begin{itemize}
    \item \textbf{Tamper-evident audit trail}: All state changes are cryptographically linked, making unauthorized modifications detectable
    \item \textbf{Leaderless coordination}: No single leader is required, reducing operational complexity and eliminating leader election overhead
    \item \textbf{Byzantine fault tolerance}: The blockchain consensus naturally handles Byzantine failures, strengthening security in a federated environment
    \item \textbf{Path to permissionless}: Blockchain-based consensus is a natural fit for eventual permissionless operation
\end{itemize}

The consensus protocol achieves finality within 2-3 seconds under normal operation, with transaction throughput sufficient for WalletConnect's write-light workload (typically $<$100 writes/second globally).

\subsection{Coordinator Decentralization}

WCN 2.0 expands the coordinator set from 5 operators in WCN 1.0 to over 20 independent operators distributed across 8+ geographic regions. Each operator runs one or more coordinator nodes that participate in blockchain consensus.

The expanded operator set improves resilience in two ways: (1) no single operator controls enough nodes to unilaterally affect consensus, and (2) geographic distribution ensures the system remains available even during regional network partitions or infrastructure failures.

\section{Performance and Benefits}

\subsection{Latency Improvements}

We evaluated WCN 2.0's latency characteristics using production traffic patterns over a 30-day period. Table~\ref{tab:latency} compares write latencies between WCN 1.0 (global replication) and WCN 2.0 (regionalized quorums).

\begin{table}[h]
\centering
\caption{Write Latency Comparison (p50/p95/p99 in ms)}
\label{tab:latency}
\begin{tabular}{@{}lccc@{}}
\toprule
Region & WCN 1.0 & WCN 2.0 & Improvement \\
\midrule
North America & 145/280/420 & 52/95/140 & 64\% (p50) \\
Europe & 165/310/485 & 58/105/165 & 65\% (p50) \\
Asia-Pacific & 180/340/520 & 65/115/180 & 64\% (p50) \\
Cross-region & 220/410/650 & 215/405/640 & 2\% (p50) \\
\bottomrule
\end{tabular}
\end{table}

Regional writes show consistent improvement of 60-65\% in median latency across all regions. Cross-region writes (requiring global quorum) show minimal degradation, validating that the promotion mechanism correctly identifies and handles global state.

Read latency remains unchanged from WCN 1.0, as both systems serve reads from local database replicas. The 99th percentile read latency averages 15ms across all regions.

\subsection{Operational Benefits}

The coordinator-database separation has significantly improved operational agility:

\begin{itemize}
    \item \textbf{Rolling upgrades}: Database upgrades can be performed with zero downtime by rotating replicas out of service. Coordinator upgrades benefit from the blockchain's fault tolerance, allowing gradual rollout across operators.
    \item \textbf{Deployment frequency}: Average deployment frequency increased from 2-3 per month in WCN 1.0 to 8-12 per month in WCN 2.0, reflecting reduced deployment risk.
    \item \textbf{Incident recovery}: Mean time to recovery (MTTR) for database failures decreased from 8 minutes to 2 minutes due to faster failover to healthy replicas.
\end{itemize}

\subsection{Security and Resilience}

WCN 2.0's decentralized architecture provides several security advantages:

\textbf{Byzantine fault tolerance}: The blockchain consensus tolerates up to $\frac{n-1}{3}$ Byzantine failures among coordinator nodes. With 20+ operators, the system can withstand 6+ malicious or faulty coordinators.

\textbf{Audit trail}: All state transitions are recorded in the blockchain with cryptographic proofs. This provides complete traceability for debugging and security analysis.

\textbf{Geographic resilience}: With 8+ regions, the system remains operational during complete regional failures. Testing confirmed availability during simulated loss of any two regions simultaneously.

\textbf{No single point of failure}: Unlike WCN 1.0's Raft leader, WCN 2.0 has no privileged nodes. Coordinator failures are handled transparently by the blockchain consensus.

\section{Hardfork Migration Strategy}

Migration from WCN 1.0 to WCN 2.0 requires careful coordination to maintain service availability during the transition. The migration follows a phased approach:

\textbf{Phase 1 - Parallel Operation}: WCN 2.0 coordinators are deployed alongside WCN 1.0 infrastructure. Both systems operate independently with shadow writes to WCN 2.0 to build up state and validate correctness.

\textbf{Phase 2 - Progressive Traffic Shift}: Client traffic is gradually shifted from WCN 1.0 to WCN 2.0 in regional increments (e.g., 5\%, 25\%, 50\%, 100\%). Monitoring validates latency improvements and error rates at each stage.

\textbf{Phase 3 - State Reconciliation}: Any state divergence between systems is identified and reconciled. The blockchain's audit trail facilitates debugging of discrepancies.

\textbf{Phase 4 - Cutover}: Once WCN 2.0 handles 100\% of traffic successfully for a sustained period (e.g., 7 days), WCN 1.0 infrastructure is decommissioned.

The migration is designed to be reversible until Phase 4, allowing rollback if critical issues emerge.

\section{Updated Reward System}

WCN 2.0 introduces an updated reward mechanism for node operators that better aligns incentives with system health and performance. The new system incorporates several improvements over WCN 1.0:

\textbf{Performance-based rewards}: Operator rewards are adjusted based on node uptime, latency performance, and consensus participation. Operators with consistent high performance receive proportionally higher rewards.

\textbf{Token-price-adjusted rewards}: To maintain stable incentives in the face of token price volatility, rewards are dynamically adjusted based on a moving average of token price over the previous 30 days. This ensures operators receive consistent economic value regardless of short-term price fluctuations.

\textbf{Geographic diversity incentives}: Additional rewards are allocated to operators in underserved regions to encourage global distribution. This strengthens network resilience and reduces latency for users worldwide.

The specific reward formula and distribution parameters are subject to ongoing governance and may be adjusted based on operational experience and community feedback.

\section{Future Work}

Several directions for future development have been identified:

\subsection{Autoscaling Storage Layer}

The current coordinator-database separation enables elastic scaling of storage, but autoscaling is not yet implemented. Future work will develop policies for automatic database replica scaling based on load metrics, further improving operational efficiency.

\subsection{Permissionless Operation}

While WCN 2.0's blockchain consensus is compatible with permissionless operation, the current deployment remains permissioned with vetted operators. The path to permissionless operation requires:
\begin{itemize}
    \item Stake-based sybil resistance mechanisms
    \item Reputation systems to identify and penalize malicious actors
    \item Dynamic operator sets that can adapt as nodes join and leave
    \item Enhanced economic security models
\end{itemize}

\subsection{Ecosystem Use Cases}

Beyond WalletConnect Protocol and Smart Sessions, WCN 2.0's architecture may serve other decentralized applications requiring low-latency, verifiable state coordination. Potential use cases include:
\begin{itemize}
    \item Decentralized identity systems
    \item Cross-chain messaging protocols
    \item Real-time collaborative applications
    \item IoT device coordination
\end{itemize}

\subsection{Governance Updates}

The reward system and operational parameters will evolve based on community governance. Future governance mechanisms may include:
\begin{itemize}
    \item On-chain voting for parameter adjustments
    \item Operator councils for technical decision-making
    \item Transparent metrics dashboards for community oversight
\end{itemize}

\section{Conclusion}

WalletConnect Network 2.0 represents a significant evolution in distributed key-value store design, introducing three major innovations: regionalized quorums for latency optimization, blockchain-based consensus for verifiable coordination, and coordinator-database separation for operational agility.

Production deployment demonstrates 60-65\% latency reduction for regional workloads while maintaining strong consistency guarantees. The expanded operator set across 20+ independent entities and 8+ geographic regions significantly improves resilience and advances WalletConnect's decentralization goals.

By replacing traditional consensus mechanisms with blockchain-based coordination, WCN 2.0 establishes a new architectural pattern for distributed systems that require both high performance and verifiable state management. The system serves as foundational infrastructure for WalletConnect's mission to build fully permissionless, globally accessible decentralized applications.

As the first production distributed key-value store to leverage blockchain for consensus, WCN 2.0 demonstrates that blockchain technology extends beyond cryptocurrency applications to solve fundamental challenges in distributed systems coordination.

\section*{Acknowledgments}

We thank the WalletConnect community and the 20+ independent node operators who participate in running WCN 2.0. Special thanks to the engineering teams who contributed to the design, implementation, and testing of this system.

% Bibliography
\bibliographystyle{plain}
\bibliography{references}

\end{document}
